\section{\method}
\label{sec:method}

\method~converts local student-teacher compatibility into a position-wise weight on OPD rewards.
The method is designed as a drop-in modifier for OPD: it does not alter how the teacher reward is computed, but changes how much each visited position should trust that reward.
It has three parts: a per-position compatibility metric, a cumulative reliability weight, and an optional dynamic response-length controller.

\subsection{Compatibility Metric}

At position $\tau$, our goal is to decide whether the teacher still provides locally exploitable supervision on the student's current trajectory.
The student and teacher are therefore evaluated on the same student-generated
prefix $s_{\tau}=(x,y_{<\tau})$.

We first restrict the comparison to each model's high-confidence candidate
tokens. Let
\begin{equation}
    \mathcal{K}^{\mathrm{S}}_{\tau}
    =
    \operatorname{TopK}\!\left(\student(\cdot\mid s_{\tau}),k\right),
    \qquad
    \mathcal{K}^{\mathrm{T}}_{\tau}
    =
    \operatorname{TopK}\!\left(\teacher(\cdot\mid s_{\tau}),k\right)
\end{equation}
be the returned top-$k$ candidate sets.
The top-$k$ size is an experimental hyperparameter.
Then, we measure local compatibility with the overlap ratio
\begin{equation}
    \overlap
    =
    \frac{
    \left|\mathcal{K}^{\mathrm{S}}_{\tau}\cap \mathcal{K}^{\mathrm{T}}_{\tau}\right|
    }{k}.
    \label{eq:local-overlap}
\end{equation}
Low overlap means that the teacher's dense reward is being computed under prefixes where its preferred next-token support has separated from the student's support. Accordingly, we define the prefix-drift event as
\begin{equation}
    B_{\tau}
    =
    \mathbf{1}\left[\overlap < \gamma\right],
    \label{eq:bad-event}
\end{equation}
where $\gamma$ is the compatibility threshold.
These drift events are later accumulated to attenuate subsequent OPD rewards and define the effective response length for dynamic rollout control.

This metric follows the OPD mechanism analysis directly: successful OPD is associated with increasing high-probability overlap, while failing runs show stagnant or unstable overlap on student-visited states \citep{li2026rethinking}.
It measures whether the student and teacher still share a plausible local action space.
Low overlap means that the teacher's dense reward is being computed in a region where its preferred next-token support has separated from the student's support.
We also design an alternative top-$p$ action-acceptance metric, its definition is given in Appendix~\ref{app:top-p-metric}.

\subsection{Cumulative Reliability and Loss Weighting}

Bad events are accumulated over the response:
\begin{equation}
    C_{\tau} = \sum_{i=1}^{\tau} B_i.
\end{equation}
The raw reliability weight is a clipped linear decay,
\begin{equation}
    R_{\tau} = \clip(1 - w_{\mathrm{drop}} C_{\tau}, 0, 1),
    \label{eq:raw-weight}
\end{equation}
where $w_{\mathrm{drop}}\geq 0$ controls the penalty per compatibility failure.
The raw weight is monotone non-increasing along each valid response.
Padding positions are forced to zero.

The actual loss weight adds a base floor:
\begin{equation}
    L_{\tau} = R_{\tau} + w_{\mathrm{base}}
    \label{eq:loss-weight}
\end{equation}
for valid response positions, and $L_{\tau}=0$ for padding.
The scaled \opd~reward is
\begin{equation}
    \widetilde{r}_{\tau,j} = L_{\tau} \cdot r_{\tau,j}.
    \label{eq:scaled-reward}
\end{equation}
The base weight has two practical roles.
First, it prevents \method~from becoming a hard switch when $R_{\tau}$ reaches zero.
Second, it lets the method trade off denoising and retention of weak teacher signal.
Concrete hyperparameter values are specified in the experimental setup.

\subsection{Dynamic Response Budget}

\method~can additionally adjust the rollout length during training.
Define the effective response length of a sample as
\begin{equation}
    E = \sum_{\tau=1}^{T} \mathbf{1}\left[R_{\tau}>\epsilon\right],
    \label{eq:effective-length}
\end{equation}
for a small tolerance $\epsilon$.
This definition deliberately uses $R_{\tau}$, not $L_{\tau}$.
Therefore a token with $R_{\tau}=0$ and $w_{\mathrm{base}}>0$ can still contribute a small \opd~reward, but it is not counted as reliable length.




% \vspace{-2mm}
\begin{algorithm}[t]
\caption{\method~with reliability-aware reward scaling and dynamic response budget}
\label{alg:prune-opd}
\begin{algorithmic}[1]
\Require prompt distribution $\mathcal{D}$, student policy $\student$, teacher policy $\teacher$, OPD reward routine, initial response budget $M_1$
% \Require compatibility, reward-weighting, and response-budget hyperparameters
\Ensure updated student policy $\student$ and response budget sequence $\{M_t\}$
\For{training step $t=1,2,\ldots$}
    \State Sample prompts $\{x^{(b)}\}_{b=1}^{B}\sim\mathcal{D}$ and generate student rollouts $\{y^{(b)}_{1:T_b}\}_{b=1}^{B}$ with maximum length $M_t$.
    \State Compute raw OPD reward tensors $r^{(b)}_{\tau,j}$ on the student-visited prefixes.
    \For{each rollout $b=1,\ldots,B$}
        \State Compute the overlap-ratio sequence on student-visited prefixes using Eq.~\eqref{eq:local-overlap}.
        \State Mark low-compatibility positions and form the cumulative reliability weights using Eqs.~\eqref{eq:bad-event}--\eqref{eq:loss-weight}.
        \State Scale OPD rewards using Eq.~\eqref{eq:scaled-reward} and compute the effective reliable length $E^{(b)}$ using Eq.~\eqref{eq:effective-length}.
    \EndFor
    \State Update $\theta$ with the OPD objective using scaled rewards $\widetilde{r}^{(b)}_{\tau,j}$.
    \State $h_t\leftarrow B^{-1}\sum_{b=1}^{B}\mathbf{1}[E^{(b)}\geq M_t-m]$.
    \If{$h_t\geq \rho$}
        \State $M_{t+1}\leftarrow \min(M_t+\Delta,M_{\max})$ and reset the low-hit counter.
    \ElsIf{$h_t<\rho$ for $P$ consecutive steps}
        \State $M_{t+1}\leftarrow \max(M_t-\Delta,M_{\min})$ and reset the low-hit counter.
    \Else
        \State $M_{t+1}\leftarrow M_t$.
    \EndIf
\EndFor
\end{algorithmic}
\end{algorithm}
% \vspace{-8mm}

% \vspace{-5mm}
Let $M_t$ be the maximum response length used for rollout generation at training step $t$.
The controller monitors the batch hit ratio: the fraction of samples whose effective reliable length reaches the current limit up to a margin.
If the hit ratio is large, the limit expands by a fixed step, capped by $M_{\max}$.
If the hit ratio remains small for a patience window, the limit contracts by the same step, bounded below by $M_{\min}$.
The initial limit, bounds, update step, margin, hit-ratio threshold, and patience are specified in the experimental setup. Algorithm~\ref{alg:prune-opd} summarizes the full \method~procedure, including reliability-aware reward scaling and dynamic response-budget adaptation.
To keep the training loop readable, the position-wise compatibility and weighting rules refer directly to the definitions above.

% The implementation is intentionally narrow.
% When \method~is disabled, the original OPD or GRPO path is unchanged.
% When enabled, \method~only rescales the already-computed OPD reward tensor.
% It does not currently route zero-reliability tokens to a GRPO fallback loss; that is an explicit future direction rather than part of the present method.